\chapter{Researching a company}
\label{ch:company}

\section*{What this chapter is for}

Three things are worth establishing before you talk to anybody: what the level
means here, what the work is actually like, and what the number is. All three
are usually findable, and none of them is findable from the posting alone.

\section{The ladder is the rubric}

Where a company publishes an engineering ladder, it is the single most useful
document available to you: it is their own statement of what the level you are
applying to means, written before they met you.

Where they do not publish one, read two or three that other companies
publish.\source{progression-fyi} The wording differs and the structure does
not, and a well-written public ladder tells you what questions the level
attracts anywhere.

\begin{kitmethod}
\textbf{1. Find the row.} Match the advertised band, or the title, to a row.

\textbf{2. Read the row above and below.} The differences between three
adjacent rows tell you more than the row itself, because they name the
transitions, and an interview is a decision about which side of a transition
you are on.

\textbf{3. Extract the verbs.} As in Chapter~\ref{ch:segment}: a ladder written
in terms of \emph{unscoped}, \emph{across domains}, \emph{influences} is
telling you which stories to choose.

\textbf{4. Notice what is not there.} A ladder with no mention of mentoring at
your level is a different job from one where it is the defining criterion, and
you cannot tell which without looking.
\end{kitmethod}

\section{The band}

Where a range is published, it is a fact about the level, not an opening
position. Two things are worth establishing about it, and both are usually
answerable in the screen.

\textbf{What the range is composed of.} A published figure may be base only,
or base plus a variable component, or include equity at a notional valuation.
Those are three different offers with the same headline. Ask which.

\textbf{Whether the range spans more than one level.} It frequently does. If
so, the top of the range is a different role from the bottom, with a different
interview, and finding out which one you are being considered for is a fair
question.

\begin{kitnote}
This book carries no salary figures and gives no advice about what any level is
worth. Those numbers are regional, they date within two quarters, and the
versions that circulate are almost all unsourced \dash{} several are named and
refused in Appendix~\ref{app:refused}. What is durable is the method: establish
the composition, establish the level span, and record what you find in
\code{templates/company-brief.md} so the next process starts from data rather
than from memory.
\end{kitnote}

\section{The work model, believed rather than hoped}

Chapter~\ref{ch:fit} covered this as a go/no-go. As research, one thing is worth
adding: check what the company's \emph{other} current postings say. A single
posting can be written loosely; a set of twenty postings where every
engineering role carries the same work-model field is a policy, and policies are
not negotiated in a screening call.

\section{Culture signals, and how much weight to put on them}

Public review sites, engineering blogs, conference talks, and the company's own
writing about how it works. All of it is evidence and all of it is biased in
knowable directions.

\begin{kittable}{@{}lXl@{}}
\textbf{Source} & \textbf{Biased by} & \textbf{Useful for} \\
\midrule
Review sites & who chooses to write a review & repeated specific complaints \\
Engineering blog & recruiting intent & what they are proud of, and what they measure \\
Conference talks & the speaker's own work & the shape of their systems \\
Their own values page & aspiration & the vocabulary the interview is scored in \\
\end{kittable}

\judgement{The useful signal on a review site is not the average. It is a
specific complaint repeated by several people about a thing that would matter
to you, particularly when it comes from the function you would be joining
rather than from the company as a whole. A low average driven by a sales
organisation says little about an engineering team, and people read it as
though it does.}

The values page deserves more attention than it usually gets, for an unsentimental
reason: where behavioural rounds are scored against stated values, that page is
the rubric. Reading it is not endorsement, it is preparation.

\begin{drill}
Fill \code{templates/company-brief.md} for the role in front of you: the ladder
row and the two either side, the band composition and level span, the work
model across their other postings, and three culture signals with their
biases named. Note every question the research could not answer \dash{} that
list is your agenda for the screening call.
\end{drill}

\section*{What this chapter does not claim}

It does not claim published ladders are honest descriptions of how promotion
works. They are statements of intent, and the gap between a ladder and its
application is a real phenomenon. They are useful because they are what the
interview will be scored against, which is a different claim from being true.

It offers no figures, no benchmarks and no ranges, deliberately, for the reason
the note above gives.
